\documentclass[11pt]{article}

% --- 1. Packages (Essential for ArXiv quality) ---
\usepackage[utf8]{inputenc}
\usepackage[T1]{fontenc}
\usepackage{amsmath,amssymb}
\usepackage{graphicx}
\usepackage{booktabs}
\usepackage{hyperref}
\usepackage{xcolor}
\usepackage{multirow}
\usepackage[numbers]{natbib}

% --- 2. Document Metadata ---
\title{TrucoBench: Evaluating LLM Strategic Reasoning Through Truco Regional Variants}     
\author{William Manzoli}
\date{June 2026}

% --- 3. The Document ---
\begin{document}

\maketitle

% --- ABSTRACT ---
\begin{abstract}
We introduce \textbf{TrucoBench}, the first benchmark for evaluating large language models (LLMs) on Truco---a popular Latin American card game featuring imperfect information, bluffing, and a unique nested escalation mechanic. We support two primary regional variants: \textbf{Truco Paulista} and \textbf{Truco Mineiro}. Unlike poker-based benchmarks, Truco's escalation system creates a distinct strategic axis where players must reason about stake-raising, opponent modeling, and fold-or-fight decisions at multiple levels. We formalize Truco as a Partially Observable Stochastic Game (POSG) and introduce the \textbf{LLMWiki} prompting technique---an information-dense format that utilizes dynamic card-strength tables to reduce cognitive load. Our results reveal that while models excel at bluffing, they struggle with hidden-information logic, a gap partially mitigated by structured domain-knowledge injection. TrucoBench and all evaluation code are open-source at \url{https://github.com/ManzoliW/trucobench}.
\end{abstract}

% --- INTRODUCTION ---
\section{Introduction}

Large language models have demonstrated remarkable capabilities across a wide range of tasks. A growing body of work evaluates LLMs on strategic games---particularly poker---as a testbed for reasoning under uncertainty and deception. However, the focus on poker leaves a gap: many culturally significant card games with distinct strategic properties remain unexplored as LLM benchmarks.

\textbf{Truco} is one of the most popular card games in Brazil and across Latin America. While it shares imperfect information with poker, it introduces a fundamentally different strategic axis: a \emph{nested escalation mechanic} where players can repeatedly raise the stakes of each round through a chain of calls (e.g., truco $\to$ seis $\to$ nove $\to$ doze). Each escalation forces the opponent into a three-way decision---accept, counter-raise, or fold---creating a rich space for bluffing, meta-reasoning, and risk assessment that is structurally distinct from poker's betting rounds. We support both the \textbf{Paulista} (variable manilhas) and \textbf{Mineiro} (fixed manilhas) variants, representing the two most common rule sets in Brazil.

We argue that Truco offers a complementary evaluation axis to existing game benchmarks:
\begin{itemize}
    \item \textbf{Escalation as nested bluffing.} Unlike poker's continuous bet sizing, Truco's discrete escalation chain creates clear decision points for commitment and credibility.
    \item \textbf{LLMWiki Prompting.} We introduce a specialized prompting format that addresses LLM reasoning failures by providing explicit, dynamic ``cheat sheets'' of card strengths, isolating strategic decision-making from basic rank calculation.
    \item \textbf{Cultural diversity in AI evaluation.} Benchmarking LLMs on games from diverse cultures helps identify potential biases in strategic reasoning.    
\end{itemize}

% --- RESULTS ---
\section{Results: Diagnostic Evaluation (Preprint v2.3)}

The diagnostic evaluation serves as a luck-independent measure of strategic ``intuition.'' We evaluated eight representative models against a curated set of 9 scenarios covering core Truco mechanics, leveraging the LLMWiki ``cheat sheet'' format to isolate strategic logic from rank calculation.

\begin{table}[ht]
\centering
\caption{Deterministic Diagnostic Accuracy (T=0.0) Across Strategic Categories (\%)}
\label{tab:diagnostics}
\begin{tabular}{@{}lccccc@{}}
\toprule
Model & Overall & Bluff & Defense & Logic & Escalation \\ 
\midrule
gemini-2.5-pro & 80.0 & 100.0 & 40.0 & 100.0 & 100.0 \\ 
gpt-4o & 61.1 & 100.0 & 100.0 & 100.0 & 60.0 \\ 
claude-sonnet-4.6 & 71.1 & 100.0 & 40.0 & 100.0 & 100.0 \\
kimi-k2.5 & 70.0 & 100.0 & 100.0 & 100.0 & 100.0 \\
deepseek-r1 & 55.6 & 100.0 & 100.0 & 100.0 & 100.0 \\
gemini-2.5-flash & 58.9 & 100.0 & 40.0 & 100.0 & 100.0 \\
gpt-4o-mini & 53.3 & 30.0 & 40.0 & 100.0 & 100.0 \\
claude-haiku-4.5 & 52.2 & 100.0 & 100.0 & 100.0 & 60.0 \\
\bottomrule
\end{tabular}
\end{table}

The results in Table~\ref{tab:diagnostics} (Deterministic $T=0.0$ configuration) reveal a significant ``Knowing-Doing Gap.'' While frontier models exhibit near-perfect accuracy in psychological categories like Bluffing, they struggle with hidden-information Logic and defensive optimization. Notably, \textbf{Gemini 2.5 Pro} and \textbf{Claude Sonnet 4.6} lead the deterministic leaderboard, followed closely by \textbf{Kimi-k2.5}.

\subsection{Ablation Study: The Inverse Information Effect}

To evaluate the impact of the \textbf{LLMWiki} technique, we conducted a deterministic ablation study across 8 models using an XML-structured format.

\begin{table}[ht]
\centering
\caption{Ablation Study (T=0.0): Standard vs. LLMWiki Prompting (Overall Accuracy \%)}
\label{tab:ablation}
\begin{tabular}{lccc}
\toprule
Model & Standard (v1) & LLMWiki (v2) & Delta \\ 
\midrule
gpt-4o & 82.2 & 61.1 & \textcolor{red}{-21.1} \\ 
gemini-2.5-pro & 74.4 & 80.0 & +\textbf{5.6} \\ 
claude-sonnet-4.6 & 80.0 & 71.1 & \textcolor{red}{-8.9} \\ 
deepseek-r1 & 66.7 & 55.6 & \textcolor{red}{-11.1} \\ 
gpt-4o-mini & 75.6 & 53.3 & \textcolor{red}{-22.2} \\ 
gemini-2.5-flash & 80.0 & 58.9 & \textcolor{red}{-21.1} \\ 
claude-haiku-4.5 & 51.1 & 52.2 & +\textbf{1.1} \\ 
kimi-k2.5 & 61.1 & 70.0 & +\textbf{8.9} \\ 
\bottomrule
\end{tabular}
\end{table}

As shown in Table~\ref{tab:ablation}, we observe a persistent \textbf{``Inverse Information Effect''} across most OpenAI and reasoning-heavy models (e.g., \textbf{GPT-4o} $-21.1\%$, \textbf{DeepSeek-R1} $-11.1\%$). In these architectures, providing explicit card-strength ``cheat sheets'' likely acted as a cognitive distraction, interfering with the models' established internal heuristics.

However, \textbf{Kimi-k2.5} (+\textbf{8.9}\%) and \textbf{Gemini 2.5 Pro} (+\textbf{5.6}\%) emerged as significant exceptions. These models appear uniquely capable of utilizing explicit XML-structured reference tables to correct or augment their internal strategic logic. Preprint v2.3 thus highlights that the utility of prompt-level domain injection is highly architecture-dependent.

\section{Conclusion}
TrucoBench provides a novel lens through which to evaluate LLM reasoning. The introduction of regional variants and LLMWiki prompting establishes a robust foundation for evaluating strategic intuition beyond standard Western benchmarks. Future work will expand to 4-player team dynamics and partner signaling.

\end{document}
